\documentclass[11pt]{article}
\usepackage[a4paper, top=18mm, bottom=18mm, left=20mm, right=20mm]{geometry}
\usepackage[T2A]{fontenc}
\usepackage[utf8]{inputenc}
\usepackage[ukrainian]{babel}
\usepackage{graphicx}
\usepackage[export]{adjustbox}
\usepackage{amsmath}
\usepackage{amssymb}
\graphicspath{{/app/Quiz/Scripts/2023/ProgAll/15BinTree/}{/app/Quiz/Scripts/2020/Mod2020_4/BinTree/}}
\setlength{\parindent}{0pt}
\setlength{\parskip}{4pt}
\pagestyle{empty}
\begin{document}
%begin
{\large\textbf{Контрольна робота}}

\textbf{Варіант \No\,1}\hfill \textit{ПІБ:}\,\underline{\hspace{60mm}}
\vspace{4mm}

%beginitem
1. Що буде виведено за виконання фрагменту коду?

%code
a = 4
b = 5
c = 2

def g(b):
global c    a-=2
    b=1
    c=4
    return a+b+c

a, b, c = 9, 8, 7
try:
    print(g(b), a, b, c, sep=';')
except:
    print('Z', a, b, c, sep=';')
%code

%enditem
%beginitem
2. %goodpath:Scripts/2018/Mod2018_1/03-good.txt
Відмітити все, що є коректним оператором С++:
( 1 ) cout<<(3=!5);

( 2 ) x=int(4);

( 3 ) if ('x'<>'y') a=6;

( 4 ) x+=y;

%enditem
%beginitem
3. Що буде виведено за виконання фрагменту коду?

%code
cout << 8 / 8 / 8 / 8;
%code

%enditem
%beginitem
4. 
Змінні $next$ та $prev$ вузла списку позначають наступний та попередній за ним вузол відповідно. Починаючи з голови, у список записано послідовні цілі числа від $-66$ до $19$. Нехай змінна $p$ позначає вузол, що зберігає значення $9$.
Яке значення зберігається у вузлі $p.next.next.next.next.next$ ?

%enditem
%beginitem
5. Що буде виведено за виконання фрагменту коду?

%code
def g(a,b):
    c=14
    if b:
        c=4
    elif b==0:
         return 5
    else: 
        c=2
    return c

print(g(9,9))
%code

%enditem
%beginitem
6. 

def f(n):
    print("f", end="");
    return n==1

print(f(9) or f(-3))

%enditem
%beginitem
7. Що буде виведено за виконання фрагменту коду?
%code
a,b,c=9,6,8
def g(a,b=7,c=6):
    print(a,b,c,end=" ")

g(a=4,5,b=2)
print(a,b,c)
%code

%enditem
%beginitem
8. Що буде виведено за виконання фрагменту коду?

%code
print('R', end=' ')
for e in range(8,8,3):
    if e>=7:
        continue
        print(e, end=' ')
        e=8
    if e>=7:
        break
    else:
        print(e, end=' ')
    print(e, end=' ')
%code

%enditem
%beginitem
9. 
Рядки журналу через двокрапку містять ім'я користувача (ідентифікатор), час події,тип події, код події, наприклад
\begin{verbatim}
s="username : 2022-05-05 13:26 : ERROR : 404"
\end{verbatim}
у коді 
\begin{verbatim}
res = re.fullmatch('???', s) 
s = ''
code_str = ??? 
\end{verbatim}
чим треба замінити кожен з ???, щоб значенням code\_str був код події? 



%enditem
%beginitem
10. Що буде виведено за виконання фрагменту коду?

%code
print('R', end=' ')
e = (0,1,2,3,4,5,6,7,8,9)
print(e[13:13:-2])
%code


%enditem
%beginitem
11. Що буде виведено за виконання фрагменту коду?

%code
print('R', end=' ')
m=4
while m<=10:
    m+=2
print(m)
%code


%enditem
%beginitem
12. 

a,b,c=4,5,2
def g(b):
global c    a-=2
    b=1
    c=4
    return a+b+c

a,b,c=9,8,7
print(g(b),a,b,c)

%enditem
%beginitem
13. Що буде виведено за виконання фрагменту коду?

%code
#include <iostream>
using namespace std;

int a = 4, b = 5, c = 2;

int g(int &b){
int c;    a = 2;
    b = 1;
    c = 4;
    return a + b + c;
}

int main(){
    inta = 9;
    intb = 8;
    intc = 7;
    cout << g(b) << ':';
    cout << a << ':' << b << ':' << c;
    return 0; 
}
%code

%enditem
%beginitem
14. 
def f(arg, n):
    n = 4
    tmp = arg
    tmp[:-5] = 'bd'
    return len(tmp)>3

try:
    e = ['0','1','2','3','4']
    n = 5
    f(e, n)
    print(n, end=';')
    for el in e:
        print(el, end=';')
except:
    print('Z')

%enditem
%beginitem
15. 

print('R', end=' ')
e=(0,1,2,3,4,5,6,7,8,9)
print(e[13:13:-2])


%enditem
%beginitem
16. 
Записати ВИРАЗ, значенням якого є словник, ключами якого є цілі числа від 15 до 2018 (включно), що діляться на 7, а ключам відповідають їх куби 
%enditem
%beginitem
17. 

def g(a,b):
    c=14
    if b:
        c=4
    elif b==0:
         return 5
    else: 
        c=2
    return c

print(g(9,9))
%enditem
%beginitem
18. 
Значенням змінної \kw{a} є цілочисловий масив типу $numpy.ndarray$ розмірів $4{\times}4$. На скільки збільшиться сума елементів масиву \kw{a} після виконання
\begin{verbatim}
a.T[1] += 2
\end{verbatim}


Якщо код некоректний, то в якості відповіді зазначити ERROR.



%enditem
%beginitem
19. Що буде виведено за виконання фрагменту коду?
%code
if (4 == 4)
    cout << "e";
else
    cout << "f";
%code


%enditem
%beginitem
20. Що буде виведено за виконання фрагменту коду?

%code
print('R', end=' ')
e=['0','1','2','3','4']
e.extend([1,2])
print(e)
%code


%enditem
%beginitem
21. Що буде виведено за виконання фрагменту коду?
%code
#include <iostream>
using namespace std;

int f(int &x, int &y){
    x = 6;
    y-= 7;
    return y;
}

int main(){
    int a = 9, b = 3;
    b = f(b, b);
    cout << a << ":" << b <<':';
    {
        int a = 5, b = 8;
        cout << ((b>=4) || ((a-=2) >= 4)) << ':';
        cout << a << ":" << b << ':';
    }
    {
        inta = 1;
        cout << a << ':';
    }
    cout << a << ":" << b << endl;
    return 0;
}
%code


%enditem
%beginitem
22. 
try:
    res = "6.0j"=="False"
    if isinstance(res, bool):
        print(f'bool;{res}')
    elif isinstance(res, int):
        print(f'int;{res}')
    elif isinstance(res, float):
        print(f'float;{res:.2f}')
    else:
        print('???')
except:
    print('Z')

%enditem
%beginitem
23. Що буде виведено за виконання фрагменту коду?

%code
#include <iostream>

bool f(int n){
    std::cout<<"f";
    return n==1;
}

int main(){
    std::cout<<(f(9) || f(-3));
    return 0;
}
%code

%enditem
%beginitem
24. Що буде виведено за виконання фрагменту коду?

%code
for e in range(8,8+3,3):
    if e>=7:
        pass
    print(e, end=' ')
    e=8
    if e>=7:
        break
else:
    print(e, end=' ')
print(e, end=' ')
%code

%enditem
%beginitem
25. Що буде виведено за виконання фрагменту коду?

%code
m=10
while m>=4:
    m+=-2
print(m, end=' ')
%code


%enditem
%beginitem
26. Що буде виведено за виконання фрагменту коду?
%code

def f(n):
    print("f", end="");
    return n==1

print(f(9) or f(-3))
%code

%enditem
%beginitem
27. 
not6==False or not7!=6.0 or 7.0>4
%enditem
%beginitem
28. Що буде виведено за виконання фрагменту коду?

%code
print('R', end=' ')
m=10
while m>=4:
    m+=-2
print(m)
%code


%enditem
%beginitem
29. 

x,y,z=4,5,2
y,z,z,y=y,x,9,z
print(x,y,z)

%enditem
%beginitem
30. 

def g(d):
    y=14
    if d: 
        y=4
    elif d==0:
         return 5
    else:
         y=2
    return y

print(g(9))

%enditem
%beginitem
31. Що буде виведено за виконання фрагменту коду?

%code
try:
    m = 10
    while m>=4:
        m += -2
    print(m, end=';')
except:
    print('Z')

%code

%enditem
%beginitem
32. 
Написати програму, що запитує у користувача значення дійсних параметрів $x$ та $y$, обчислює для них значення виразу 
$$\frac{\log_{10} x - 53}{(\arccos x - 0.6) (y - 97)}$$ 
та повiдомляє введенi данi та результат обчислення.
 Оцінюється правильність та якість коду.


%enditem
%beginitem
33. 
a = 4
b = 5
c = 2

def g(b):
global c    a = 2
    b = 1
    c = 4
    return a + b + c

a = 9
b = 8
c = 7

try:
    print(g(b), a, b, c, sep=';')
except:
    print('Z', a, b, c, sep=';')

%enditem
%beginitem
34. Що буде виведено за виконання фрагменту коду?

%code
e = ('a','b','c','d','e',0,1,2)
print(e[-8])
%code


%enditem
%beginitem
35. 
try:
    print(4, end=';')
    print(5//0.0, end=';')
    print(2, end=';')
except ValueError:
    print(9, end=';')
except Exception:
    print(8, end=';')
except:
    print('Z', end=';')
else:
    print(7, end=';')
finally:
    print(3)

%enditem
%beginitem
36. 

print('R', end=' ')
m=10
while m>=4:
    m+=-2
print(m)


%enditem
%beginitem
37. 
def f(n):
    print('f', end='')
    return n==1

try:
    print(f(9) or f(-3))
except:
    print('ERR')

%enditem
%beginitem
38. 
$a$ є цілочисловим масивом типу $numpy.ndarray$ розмірів $5{\times}5$. Сума елементів масиву $a$ дорівнює 48. Чому дорівнює сума елементів масиву $a$ після виконання 
\begin{verbatim}
a -= 48
\end{verbatim}


Якщо код некоректний, то в якості відповіді зазначити ERROR.



%enditem
%beginitem
39. Що буде виведено за виконання фрагменту коду?

%code
#include <iostream>
using namespace std;

int g(int d){
    int y = 14;
    if (d) 
        y = 4;
    else if (d == 0)
         return 5;
    else
         y = 2;
    return y;
}

int main(){
    cout << g(9);
    return 0;
}
%code


%enditem
%beginitem
40. 
Задано тип вузла списку
struct Node \{int n; Node *prev, *next;\};
У списку зберігаються послідовні цілі числа від -66 до 19. Вказівник p встановлено на вузол, в якому зберігається число 9.
Чому дорівнює значення p->next->next->next->next->next->n ?
%enditem
%beginitem
41. Що буде виведено за виконання фрагменту коду?

%code
m=4
while m<=10:
    m+=2
print(m, end=' ')
%code


%enditem
%beginitem
42. Що буде виведено за виконання фрагменту коду?

%code
m=10
while m>=4:
    m+=-2
print(m)
%code


%enditem
%beginitem
43. %goodpath:Scripts/2019/Mod2019_1/Lex/03-good.txt
Відмітити все, що є коректним оператором С++:
( 1 ) x=*2

( 2 ) x,y,z=z,z,y

( 3 ) x not is y

( 4 ) "\n" is not '\n'

%enditem
%beginitem
44. 

\begin{center}
	\adjustimage{max size={0.8\linewidth}{0.8\paperheight}}
	{../BinTree/trees_exp/168.png}
\end{center}
Указати послідовність міток вершин дерева, зображеного на рис., що утвориться в результаті обходу дерева в оберненому порядку. Мітки записувати через пробіл.\par Обхід дерева починається з кореня (на рис. це верхня вершина).
%answer:168 оберненому

%enditem
%beginitem
45. Що буде виведено за виконання фрагменту коду?

%code
a,b,c=4,5,2
def g(b):
    a-=2
    b=1
    c=4
    return a+b+c

a,b,c=9,8,7
print(g(b),a,b,c)
%code

%enditem
%beginitem
46. 

a,b,c=9,6,8
def g(a,b=7,c=6):
    print(a,b,c,end="")

g(a=4,5,b=2)
print(a,b,c)


%enditem
%beginitem
47. 

print('R', end=' ')
m=4
while m<=10:
    m+=2
print(m)


%enditem
%beginitem
48. Що буде виведено за виконання фрагменту коду?

%code
8//8%8*8
%code

%enditem
%beginitem
49. Що буде виведено за виконання фрагменту коду?

%code
print('R', end=' ')
e=(0,1,2,3,4,5,6,7,8,9)
print(e[13:13:-2])
%code


%enditem
%beginitem
50. 

print('R', end=' ')
e=['0','1','2','3','4']
e[:-5]='13'
print(e)


%enditem
%beginitem
51. 
try:
    res = not6==False or not7!=6.0 or 7.0>4
    if isinstance(res, bool):
        print(f'bool;{res}')
    elif isinstance(res, int):
        print(f'int;{res}')
    elif isinstance(res, float):
        print(f'float;{res:.2f}')
    else:
        print('???')
except:
    print('ERR')

%enditem
%beginitem
52. 
Кожен рядок текстового файлу містить по два дійсних числа, розділені білими символами.
Написати функцію, що для заданого текстового файлу будує новий текстовий файл, рядки якого крім зображень дйсних чисел ще містять результат обчислення на них функції $f$. У вихідному файлі зображення чисел мають розділятися пробілами. 
$$f(x, y) =\frac{\log_{10} x - 53}{(\arccos x - 0.6) (y - 97)}$$ 
Якість коду також оцінюється.


%enditem
%beginitem
53. Що буде виведено за виконання фрагменту коду?

%code
print('R', end=' ')
e = ('a','b','c','d','e',0,1,2)
print(e[-8])
%code


%enditem
%beginitem
54. Що буде виведено за виконання фрагменту коду?

%code
e=['0','1','2','3']
e.extend([1,2])
print(e)
%code


%enditem
%beginitem
55. Що буде виведено за виконання фрагменту коду?

%code
#include <iostream>
using namespace std;

int g(int a, int b){
    int c = 14;
    if (b)
        c = 4;
    else if (b == 0)
         return 5;
    else 
        c = 2;
    return c;
}

int main(){
    cout << g(9, 9);
    return 0;
}
%code

%enditem
%beginitem
56. 
8//8%8*8
%enditem
%beginitem
57. Що буде виведено за виконання фрагменту коду?
%Оригинал был утерян. Требует отладки!
%code
__d = 0

class D:
    __d = 1
    
    def __init__(self):
        self.__d = 2
        __d = 3
        D.__d = 4

obj = D()
D.__d = 5
try:
    print(__d, end=' ')
    print(obj.__d, end=' ')
    print(D.__d, end=' ')
    print(obj.__d, end=' ')
    print(obj._D__d)
except:
    print('ERR')
%code

%enditem
%beginitem
58. %goodpath:Scripts/2019/Mod2019_1/Lex/01-good.txt
Відмітити все, що є однією правильною лексемою:
( 1 ) ++

( 2 ) fop

( 3 ) throw

( 4 ) true

%enditem
%beginitem
59. Що буде виведено за виконання фрагменту коду?
%code
#include <iostream>
using namespace std;

int f(int &x, int &y){
    x = 6;
    y-= 7;
    return y;
}

int main(){
    int a = 9, b = 3;
    b = f(b, b);
    cout << a << ":" << b <<':';
    {
        int a = 5, b = 8;
        cout << ((b>=4) || ((a-=2) >= 4))
             << a << ":" << b << ":";
    }
    {
        inta = 1;
        cout << a << ':';
    }
    cout << a << ":" << b << endl;
    return 0;
}
%code


%enditem
%beginitem
60. Що буде виведено за виконання фрагменту коду?

%code
print(not6==False or not7!=6.0 or 7.0>4)
%code

%enditem
%beginitem
61. 
a = 4
b = 5
c = 2

def g(b):
global c    a = 2
    b = 1
    c = 4
    return a + b + c

a = 9
b = 8
c = 7

try:
    print(g(b), a, b, c, sep=';')
except:
    print('ERR', a, b, c, sep=';')

%enditem
%beginitem
62. 
Записати ВИРАЗ, значенням якого є список, що складається з цілих чисел від 15 до 2018 (включно), причому спочатку за спаданням записано всі, що діляться на 7, а потім решту за спаданням.
%enditem
%beginitem
63. %goodpath:Scripts/2018/Mod2018_1/01-good.txt
Відмітити все, що є однією правильною лексемою:
( 1 ) **

( 2 ) /

( 3 ) (1)

( 4 ) "1.1.1"

%enditem
%beginitem
64. Що буде виведено за виконання фрагменту коду?

%code
cout << (8 / 8 / 8 / 8);
%code

%enditem
%beginitem
65. Що буде виведено за виконання фрагменту коду?
%code
if (4 == 4)
    cout << "e";
else
    cout << "e";
%code


%enditem
%beginitem
66. Що буде виведено за виконання фрагменту коду?

%code
try:
    m = 4
    while m<=10:
        m += 2
    print(m, end=';')
except:
    print('Z')

%code

%enditem
%beginitem
67. Що буде виведено за виконання фрагменту коду?

%code
#include <iostream>
using namespace std;

bool f(int n){
    cout<<"f";
    return n==1;
}

int main(){
    cout<<(f(9) || f(-3));
    return 0;
}
%code

%enditem
%beginitem
68. Що буде виведено за виконання фрагменту коду?
code
d = 0
j = 1

class D:
    j = 2
    
    def __init__(self):
global j        self.j = 3
        j = 4
        D.j = 5

obj = D()
print(d, j, D.j, obj.j, sep=' ')
code

%enditem
%beginitem
69. Що буде виведено за виконання фрагменту коду?

%code
#include <iostream>

int g(int d){
    int y = 14;
    if (d) 
        y = 4;
    else if (d == 0)
         return 5;
    else
         y = 2;
    return y;
}

int main(){
    std::cout << g(9);
    return 0;
}
%code


%enditem
%beginitem
70. Що буде виведено за виконання фрагменту коду?

%code
cout << (9 / 9 / 9 / 9);
%code

%enditem
%beginitem
71. Що буде виведено за виконання фрагменту коду?
%code

x,y,z=4,5,2
y,z,z,y=y,x,9,z
print(x,y,z)
%code

%enditem
%beginitem
72. 
try:
    res = 8//8%8*8
    if isinstance(res, bool):
        print(f'bool;{res}')
    elif isinstance(res, int):
        print(f'int;{res}')
    elif isinstance(res, float):
        print(f'float;{res:.2f}')
    else:
        print('???')
except:
    print('ERR')

%enditem
%beginitem
73. Що буде виведено за виконання фрагменту коду?

%code
cout << (!6==false || !7!= 6.0 || 7.0>4);
%code

%enditem
%beginitem
74. 
Записати ВИРАЗ, значенням якого є множина, що складається з квадратних коренівцілих чисел від 15 до 2018 (включно), що діляться на 7.
%enditem
%beginitem
75. 

print('R', end=' ')
e=['0','1','2','3','4']
e.extend([1,2])
print(e)


%enditem
%beginitem
76. 
try:
    res = "6.0j"=="False"
    if isinstance(res, bool):
        print(f'bool;{res}')
    elif isinstance(res, int):
        print(f'int;{res}')
    elif isinstance(res, float):
        print(f'float;{res:.2f}')
    else:
        print('???')
except:
    print('ERR')

%enditem
%beginitem
77. 
try:
    print(4, end=';')
    print(5//0.0, end=';')
    print(2, end=';')
except ValueError:
    print(9, end=';')
except Exception:
    print(8, end=';')
except:
    print('ERR', end=';')
else:
    print(7, end=';')
finally:
    print(3)

%enditem
%beginitem
78. Що буде виведено за виконання фрагменту коду?

%code
not6==False or not7!=6.0 or 7.0>4
%code

%enditem
%beginitem
79. 
4**-1.0*True
%enditem
%beginitem
80. 
try:
    res = not6==False or not7!=6.0 or 7.0>4
    if isinstance(res, bool):
        print(f'bool;{res}')
    elif isinstance(res, int):
        print(f'int;{res}')
    elif isinstance(res, float):
        print(f'float;{res:.2f}')
    else:
        print('???')
except:
    print('Z')

%enditem
%beginitem
81. 

\begin{center}
	\adjustimage{max size={0.8\linewidth}{0.8\paperheight}}
	{../15BinTree/trees_exp/168.png}
\end{center}
Указати послідовність міток вершин дерева, зображеного на рис., що утвориться в результаті обходу дерева в оберненому порядку. Мітки записувати через пробіл.\par Алгоритм починає роботу з кореня (на рис. це верхня вершина).
%answer:168 оберненому

%enditem
%beginitem
82. Що буде виведено за виконання фрагменту коду?

%code
4**-1.0*True
%code

%enditem
%beginitem
83. 
try:
    res = 4**-1.0*True
    if isinstance(res, bool):
        print(f'bool;{res}')
    elif isinstance(res, int):
        print(f'int;{res}')
    elif isinstance(res, float):
        print(f'float;{res:.2f}')
    else:
        print('???')
except:
    print('Z')

%enditem
%beginitem
84. Що буде виведено за виконання фрагменту коду?

%code
"6.0j"=="False"
%code

%enditem
%beginitem
85. 
e=['0','1','2','3','4']; e[:-5]='13'; print(e)
%enditem
%beginitem
86. Що буде виведено за виконання фрагменту коду?

%code
print(6 != False is 7 >= 6.0)
%code

%enditem
%beginitem
87. Що буде виведено за виконання фрагменту коду?

%code
cout << 9 / 9 / 9 / 9;
%code

%enditem
%beginitem
88. Що буде виведено за виконання фрагменту коду?
%code
d = 0
j = 1

class D:
    j = 2
    
    def __init__(self):
global j        self.j = 3
        j = 4
        D.j = 5

obj = D()
try:
    print(d, end=' ')
    print(j, end=' ')
    print(D.j, end=' ')
    print(obj.j)
except:
    print('ERR')
%code

%enditem
%beginitem
89. Що буде виведено за виконання фрагменту коду?
%code
for e in range(13,13,-2):
    if e>=7:
        continue
        print(e,end=' ')
        e=8
    if e>=7:
        break
else:
    print(e, end=' ')
print(e, end=' ')
%code

%enditem
%beginitem
90. 
Змінні \kw{next} та \kw{prev} вузла списку позначають наступний та попередній за ним вузол відповідно. Починаючи з голови, у список записано послідовні цілі числа від~$-66$ до~$19$. Нехай змінна \kw{p} позначає вузол, що зберігає значення~$9$.
Яке число зберігається у вузлі \kw{p.next.next.next.next.next} ?

%enditem
%beginitem
91. Що буде виведено за виконання фрагменту коду?
%code
a, b, c = 9, 6, 8
def g(a, b=7, c=6):
    print(a, b, c, end=" ")

g(a=4, 5, b=2)
print(a, b, c)
%code

%enditem
%beginitem
92. 
Написати програму, що запитує у користувача значення дійсних параметрів $x$ та $y$, обчислює для них значення виразу 
$$\frac{\log_{10} x - 53}{(\arccos x - 0.6) (y - 97)}$$ 
та повiдомляє введенi данi та результат обчислення.


%enditem
%beginitem
93. Що буде виведено за виконання фрагменту коду?

%code
#include <iostream>

int g(int a, int b){
    int c = 14;
    if (b)
        c = 4;
    else if (b == 0)
         return 5;
    else 
        c = 2;
    return c;
}

int main(){
    std::cout << g(9, 9);
    return 0;
}
%code

%enditem
%beginitem
94. 
В умовах попередньої задачі було виконано p->next=p->prev->next;

Чому дорівнює значення p->next->next->next->next->n ?
%enditem
%beginitem
95. 
def f(n):
    print('f', end=';')
    return n==1

try:
    print(f(9) or f(-3))
except:
    print('Z')

%enditem
%beginitem
96. 
try:
    for e in range(13, 13, -2):
        if e>=7:
            continue
        print(e, end=';');e = 8
        if e>=7:
            break
    else:
        print(e,end=';')
    print(e)
except:
    print('Z')

%enditem
%beginitem
97. 

print('R', end=' ')
for e in range(13,13,-2):
    if e>=7:
        continue
        print(e, end=' ')
        e=8
    if e>=7:
        break
    else:
        print(e, end=' ')
    print(e)

%enditem
%beginitem
98. 
try:
    m = 4
    while m<=10:
        m += 2
    print(m, end=';')
except:
    print('Z')

%enditem
%beginitem
99. Що буде виведено за виконання фрагменту коду?

%code
print(8//8%8*8)
%code

%enditem
%beginitem
100. Що буде виведено за виконання фрагменту коду?

%code
print("6.0j" == "False")
%code

%enditem




































































































%end
\newpage
\end{document}


